% ============================================================
% 第 14 章  一般谱网络与 Fenchel–Nielsen 谱网络
% 第 15 章  非阿贝尔化与阿贝尔化
% 底本：Springer LNM 2386，Chapters 14–15，印刷 pp.163–180 / PDF p203–p220
% 交付：work/src/part/ch14_15.tex（片段，不含导言区）
% ============================================================

\chapter{一般谱网络与 Fenchel–Nielsen 谱网络}

% [原文 p.163]
\ChapterAbstract{本章研究更一般（退化）的谱网络，它们对许多物理应用都十分重要。这类一般网络的一个特殊类型出现在下述情形：某个 Strebel 微分的临界图的所有叶子的两端都落在分支点上。这样的谱网络称为 Fenchel–Nielsen 谱网络，因为它们与某些长度–扭转参数有关；这一构造将在 §15.6 中演示。}

\section{一般谱网络}

% [原文 p.163]
我们将把谱网络的定义作进一步的推广，以涵盖谱线相互重叠的情形，以及谱线终止于分支点而非穿孔的情形。这一层次的普遍性对第 IV 部分非常重要，但我们在此不描述它们的路径提升映射，因为一般情形下的构造在技术上更为复杂。关于路径提升映射的详细描述，我们参见 [210]。

设 $\pi\colon\overline{\Sigma}\to\overline{S}$ 是一个分歧的 $k$ 重覆叠。

\begin{definition}[14.1]
联系于 $\pi$ 的秩 $k$ 的\textbf{一般（或退化）谱网络}（general (or degenerate) spectral network）$\mathcal{W}$ 由 $\Sigma$ 上两族有限的光滑、简单、带定向的道路组成：
\[
\{p_h\colon[-1,1]\to\overline{\Sigma}\}_{h\in H_1}
\qquad\text{与}\qquad
\{p_h\colon[-1,1]\setminus\{0\}\to\overline{\Sigma}\}_{h\in H_2} ,
\]
其中 $H_1$ 与 $H_2$ 是两个不交的指标集。这些道路只能横截相交或者重叠，并且满足下列条件：
\begin{enumerate}[label=\arabic*.]
\item 对一切 $h\in H_1\cup H_2$ 与一切 $t\in(0,1]$，有 $\pi(p_h(t))=\pi(p_h(-t))$；
\item 对一切 $h\in H_1$，$\pi(p_h(0))\in B$ 是一个分支点；
\item 对一切 $h\in H_2$，下列极限存在：$p_h(+0):=\lim_{t\to+0}p_h(t)$ 与 $p_h(-0):=\lim_{t\to-0}p_h(t)$。此外 $p_h(+0)\neq p_h(-0)$，但 $\pi(p_h(+0))=\pi(p_h(-0))$；
\item 对一切 $h\in H_1\cup H_2$，$\pi(p_h(1))=\pi(p_h(-1))=s_j\in P$（对某个 $j$）是下列二者之一：
\begin{enumerate}[label=\alph*.]
\item 一个穿孔，此时相对于 $\pi^{-1}(s_j)$ 上的序有 $p_h(-1)<p_h(1)$；
% [原文 p.164]
\item 一个分歧点 $b'$，此时还存在一条经过 $b'$ 并与 $p_h$ 反向重叠的谱线 $p_{h'}$（见图 14.2）。这一对相互重叠的谱线称为\textbf{双墙}（double wall）或\textbf{双向街}（two-way street）；
\end{enumerate}
\item 任一分支点的邻域都如图 14.1 所示；
\item $\overline{\Sigma}$ 中不同的道路 $p_h$ 与 $p_{h'}$ 只能在一个点 $p_h(t)=p_{h'}(t')$ 处相遇，并且下列之一成立：
\begin{enumerate}[label=\alph*.]
\item $h,h'\in H_1$，$t=t'=0$，且 $p_h(t)=p_{h'}(t')$ 是一个分歧点；
\item $h,h'\in H_1\cup H_2$，$t,t'\in\{-1,1\}$，且 $p_h(t)=p_{h'}(t')$ 是一个穿孔；
\item $h,h'\in H_1\cup H_2$，$t\cdot t'>0$，此时我们还要求：除了与 $p_h$ 或 $p_{h'}$ 重叠的那条谱线之外，$\mathcal{W}$ 中没有任何谱线经过 $p_h(-t)$ 或 $p_{h'}(-t')$；
\item $h,h'\in H_1\cup H_2$，$t\cdot t'<0$，此时不妨设 $t<0$、$t'>0$。在这种情况下，存在一条谱线 $p_{h''}$，$h''\in H_2$，使得 $p_{h''}(+0)=p_h(-t)$ 且 $p_{h''}(-0)=p_{h'}(-t')$。此外，除了与 $p_h$、$p_{h'}$ 或 $p_{h''}$ 重叠的谱线之外，$\mathcal{W}$ 中没有任何其他谱线经过 $p_h(-t)$ 或 $p_{h'}(-t')$。
\end{enumerate}
\end{enumerate}
\end{definition}
% [Tutor Note: 原文条件 4a 引用 $\pi^{-1}(s_j)$ 上的一个序（“with respect to the order on $\pi^{-1}(s_j)$”），但前文从未定义该序；按上下文应理解为：在穿孔 $s_j$ 处取定一个小圆周（或取 $\Sigma\to S$ 在 $s_j$ 附近的局部模型），把 $\pi^{-1}(s_j)$ 的 $k$ 个点按该圆周上的循环顺序（等价地，按叶的编号）排成一个序，于是条件 4a 说的是谱线 $p_h$ 的终点 $p_h(1)$ 比起始点 $p_h(-1)$ “更靠后”，即谱线在 $\Sigma$ 上不封闭、而是从一叶绕到另一叶。核对方法：$k=2$ 且 $\pi^{-1}(s_j)=\{p_+,p_-\}$ 时该序只有两种取法，取反则 4a 的不等号同时反向，而“$p_h(-1)\neq p_h(1)$”这一实质内容不变；这说明原文所用的是一个任取的辅助序，不影响定义的实质。正文已按原文照译，未加解释性文字。]
% [Tutor Note: 条件 6 是“横截性 + 重叠相容性”的完备清单，用来排除三条以上谱线在同一内点相交的退化情形。其中 6c（$t\cdot t'>0$，两端参数同号）与 6d（$t\cdot t'<0$，异号）的区别在于两线在该点是否“同向”通过：同向时只允许另有与之重叠的谱线，异向时必须存在第三条谱线 $p_{h''}\in H_2$ 把两条半线接起来。核对方法：对一条双墙（$p_h$ 与 $p_{h'}$ 反向重叠）取交点处的参数，重叠意味着同一条几何道路被两个相反参数化，故 t 与 t' 异号，落在 6d；而 6c 处理的是两条真正不同方向的谱线恰好同向相切的情形。]

\Fig[0.55\linewidth]{figures/fig_14_1.png}{14.1}{$S$ 中一个分支点的邻域（左），以及 $\Sigma$ 中一个分歧点的邻域（右）。注意任一条谱线都可以成为双墙。}

与以往一样，我们记
\[
\mathcal{W}_\Sigma:=\bigcup_{h\in H_1\cup H_2}p_h\big((-1,1)\big)\subseteq\Sigma
\]
为 $\Sigma$ 上的谱网络，并记
\[
\mathcal{W}_S:=\pi(\mathcal{W}_\Sigma)\subseteq S
\]
为 $S$ 上的谱网络。

给定一个至少含有一条双墙的一般谱网络，我们称对每条双墙分别选取“英式”或“美式”为一个\textbf{消解}（resolution）。我们把这一选取理解为将双墙无穷小地劈开的方式，这一术语来自英国与美国的道路通行惯例。双墙的两种消解见图 14.2。当我们处理一般谱网络时，我们将始终假定已经作好了一个消解的选择。

\Fig[0.8\linewidth]{figures/fig_14_2.png}{14.2}{画在 $S$ 上的双墙的美式消解（左）与英式消解（右）。标号为 $ij$ 的谱线从 $\Sigma$ 的第 $i$ 叶走到第 $j$ 叶。注意按定义，双墙的每条谱线恰好经过两个分支点中的一个，另一个则是 $t\to\pm1$ 时的极限。}

% [原文 p.165]
\section{Fenchel–Nielsen 谱网络}

一般谱网络的定义允许谱线终止于分支点而非穿孔，只要它们是双墙即可。于是人们可以设想一个只有双墙、没有谱线终止于穿孔的谱网络。这样的谱网络甚至可以定义在闭曲面上，因为它不需要任何穿孔的存在。这样的谱网络确实存在，下面我们就来定义它们的一个特殊类型，称为 \textbf{Fenchel–Nielsen 谱网络}。这一类特殊的谱网络由 Hollands 与 Neitzke 在 [205] 中引入，之所以这样称呼，是因为这类谱网络与 §7.1 中定义的 Fenchel–Nielsen 坐标有很强的关联。这一关联将在 §15.6 中进一步研究。这些特殊的 WKB 谱网络出现在考虑 Strebel 微分（见 §8.5）的临界图之时。

设 $\overline{S}$ 是闭黎曼曲面，$P=\{s_1,\dots,s_n\}$ 是 $\overline{S}$ 上有限个标记点，$n\in\Z_{>0}$，令 $S=\overline{S}\setminus P$。设 $\omega$ 是 $\overline{S}$ 上的亚纯二次微分，在 $P$ 的各个点处有二阶极点。选取相位 $\vartheta\in\R/2\pi\Z$，使得 $\phi:=\ee^{-2\ii\vartheta}\omega$ 按定义 8.14 的意义是 Strebel 微分。这将意味着存在 $S$ 的一个裤子分解，使得 $\sqrt{\omega}$ 绕每条腿的周期具有相位 $\vartheta$。亚纯二次微分 $\omega$ 的平方根定义一个二重覆叠 $\pi\colon\Sigma_\omega\to S$，
\[
\Sigma_\omega:=\{(z,\lambda)\in S\times T_z^\ast S:\lambda^2=\omega(z)\} ,
\]
其中 $T_z^\ast S$ 是 $S$ 的余切丛在 $z$ 处的纤维。这里 $\pi(z,\lambda)=z$，而覆叠 $\pi\colon\Sigma\to S$ 的分支点是 $\omega$ 的零点。与 §13.6 中 WKB 谱网络的构造类似，我们把谱网络 $\mathcal{W}(\phi)$ 构造为临界图 $\hat{\mathcal{F}}_\phi$ 的一组叶，其中所有叶的两端都落在分支点上；于是每个叶给出网络 $\mathcal{W}(\phi)$ 的一条双墙。这样的谱网络称为 \textbf{Fenchel–Nielsen 谱网络}。

\begin{remark}[14.2]
把相位 $\vartheta$ 扰动为 $\vartheta_+=\vartheta+\epsilon$（相应地 $\vartheta_-=\vartheta-\epsilon$），其中 $\epsilon>0$ 充分小，则 $\phi$ 的临界叶分裂为两个正则叶，连接一个分支点%
% [原文 p.166]
与一个穿孔。当 $\epsilon\to0$ 时的极限便形成一条双墙，对应于 $\mathcal{W}(\phi)$ 的英式（相应地，美式）消解。
\end{remark}

$S\setminus\mathcal{W}(\phi)$ 的每个连通分支都是被闭的 $\vartheta$-轨道扫出的一个穿孔圆盘，而每个这样的分支的边界是由一条或多条\textbf{鞍连接}（saddle connection）构成的多边形。对每个这样的穿孔圆盘，Strebel 微分在穿孔 $s_j$（$j=1,\dots,n$）处的系数 $m_j\ee^{-\ii\vartheta}$，等价地，该穿孔圆盘中任一闭水平 $\vartheta$-轨道的长度，是一个正实数，它给 $S\setminus\mathcal{W}(\phi)$ 的每个连通分支打上标签。

一个 Fenchel–Nielsen 谱网络诱导出曲面 $S$ 的一个裤子分解，而在每条裤子上，该谱网络属于两种类型之一，[205] 中称它们为 \textbf{分子 I}（molecule I）与 \textbf{分子 II}（molecule II），见图 14.3。
% [Tutor Note: “每个分支都是被闭水平 $\vartheta$-轨道扫出的穿孔圆盘”可如下核对：在二阶极点 $s_j$ 处，$\phi$ 的标准局部形式为 $\phi\approx-\big(\frac{m_j}{2\pi}\big)^2\frac{\dd z^2}{z^2}$，其水平轨道是同心圆，故极点邻域内 $S\setminus\mathcal{W}(\phi)$ 的分支同胚于圆盘挖去极点；而两条临界叶（水平轨道的极限）给出该圆盘边界的两条边。于是边界多边形的每条边都是 $\phi$ 的一条鞍连接，即临界轨道的一段，其两端为分支点。]

\Fig[0.55\linewidth]{figures/fig_14_3.png}{14.3}{裤子上的两种可能的 Fenchel–Nielsen 谱网络：分子 I（上）与分子 II（下）。波浪线是分支割线，两个谱网络都采用英式消解。}

\begin{remark}[14.3]
对亚纯二次微分 $\omega$ 还存在另一种非一般的情形，此时可以构造出有趣的谱网络。即：可用同样的方式构造出一个同时含单墙与双墙的 WKB 谱网络。这样的谱网络称为\textbf{混合的}（mixed）；其实例可在 [205] 中找到，这些网络正是在该文中首次引入的。另一种有趣的%
% [原文 p.167]
谱网络情形同样由 [205] 引入，涉及所谓\textbf{收缩谱网络}（contracted spectral networks），它们是把一个 Fenchel–Nielsen 网络中的某些双墙塌缩到一起而得到的；换言之，是考虑某个圆环收缩到零尺寸时通常 Fenchel–Nielsen 谱网络的退化。
\end{remark}

% [原文 p.168]
% 本页为原书空白页（第 14 章章末的空页），无正文内容。

\chapter{非阿贝尔化与阿贝尔化}

% [原文 p.169]
\ChapterAbstract{作为迄今所引入谱网络的一个重要应用，本章给出非阿贝尔化映射与阿贝尔化映射的构造。这些映射刻画了黎曼曲面 $S$ 上秩 $n$ 丛的平坦联络，与谱曲线 $\Sigma$（$S$ 的 $n$ 叶分歧覆叠）上线丛的平坦联络之间的等价性。阿贝尔化后局部系统的单值（monodromy）正是簇 $\mathcal{X}$-坐标。因此（非）阿贝尔化的过程为在高阶泰希米勒空间的研究中使用谱网络提供了非常强的动机。特别地，我们展示如何用谱网络与阿贝尔化，在平坦联络的空间上诱导出 Fock–Goncharov 坐标或 Fenchel–Nielsen 坐标。论证谱网络与高阶泰希米勒理论之间联系的合适文献包括 [154]、[204]、[205]、[207]、[287]、[288]。}

\section{非退化谱网络的非阿贝尔化映射}

% [原文 p.169]
设 $S$ 是纤毛曲面，$\pi\colon\overline{\Sigma}\to\overline{S}$ 是带有非退化谱网络 $\mathcal{W}$ 的分歧 $k$ 重覆叠。与以往一样，$\pi$ 在 $\Sigma$ 上的分歧轨迹记作 $B'$，并记 $B:=\pi(B')\subset S$。设 $\mathcal{E}^{\mathrm{tw}}(\Sigma,\C^\times)$ 为 $\Sigma\setminus B'$ 上、在分歧点周围单值为 $-1$ 的平坦线丛所成的空间。

我们在空间 $\mathcal{E}^{\mathrm{tw}}(\Sigma,\C^\times)$ 中取一点，它给出 $\Sigma\setminus B'$ 上的一个线丛 $L\to\Sigma\setminus B'$，并带有平坦联络 $\nabla$。对 $q\in\Sigma\setminus B'$，我们用 $L_q$ 记 $L\to\Sigma\setminus B'$ 在 $q$ 上的纤维，它是一个一维的 $\C$-向量空间。

此外，我们固定 $p\in S\setminus\pi(\mathcal{W})$。我们如下构造一个表示 $\rho\colon\pi_1(S,p)\to\GL_k(\C)$。考虑如下 $k$ 维向量空间：
\[
V_p:=\bigoplus_{i=1}^{k}L_{p_i} ,
\]
% [原文 p.170]
其中 $\{p_1,\dots,p_k\}=\pi^{-1}(p)$。$V_p$ 有自然基，即 $(e_1,\dots,e_k)$，其中 $(e_i)$ 是 $L_{p_i}$ 的基，$i\in\{1,\dots,k\}$。

$S$ 上每个满足 $[\gamma]\in\pi_1(S,p)$ 的回路 $\gamma$ 都可以关于谱网络 $\mathcal{W}$ 提升到 $\Sigma$。于是我们得到一个元素 $x\in\mathbb{Z}^{\mathrm{tw}}_{B'}[\mathcal{P}(\Sigma)]$，按构造它是 $\Sigma$ 上曲线 $\gamma_r$（$r\in I$，$I$ 为有限指标集）的有限和，使得
\begin{equation}
\gamma_r(0),\gamma_r(1)\in\pi^{-1}(p) , \tag{15.1}
\end{equation}
且 $x=\sum_{r\in I}\gamma_r$。我们考虑向量 $w=\sum_{i=1}^{k}a_ie_i\in V_p$，$i\in\{1,\dots,k\}$。由于 (15.1)，沿 $\gamma$ 的平行移动可以定义为
\[
T_\gamma:=\sum_{r\in I}T_{\gamma_r}\in\bigoplus_{\tilde p,\tilde q\in\pi^{-1}(p)}L(E_{\tilde p},E_{\tilde q}) ,
\]
其中 $T_{\gamma_r}$ 是沿 $\gamma_r$ 的平行移动，由 $\Sigma$ 上的联络 $\nabla$ 给出。

对每个满足 $\gamma_r(0)=p_i$、$\gamma_r(1)=p_j$ 的 $\gamma_r$（$i,j\in\{1,\dots,k\}$），我们可以考虑 $T_{\gamma_r}e_i=t_re_j\in E_{p_j}$。我们可以按如下法则把 $T_{\gamma_r}$ 扩充为 $V_p$ 上的线性映射：
\[
T_{\gamma_r}e_l=
\begin{cases}
t_re_j, & l=i\\
0, & l\neq i .
\end{cases}
\]
我们按法则 $T_\gamma:=\sum_{r\in I}T_{\gamma_r}$ 定义线性映射 $T_\gamma\colon V_p\to V_p$。按构造，这个映射是线性的；并且由于借助谱网络 $\mathcal{W}$ 的路径提升是同伦不变的，而 $\Sigma$ 上的联络是平坦的，映射 $T_\gamma$ 只依赖于 $\gamma$ 在 $\pi_1(S,p)$ 中的同伦类。
% [Tutor Note: 式 (15.1) 之所以关键，是因为它保证每个 $\gamma_r$ 的两端都落在同一纤维 $\pi^{-1}(p)$ 上，于是 $T_{\gamma_r}$ 是 $L$ 的各纤维之间的同构，可以相加而得到 $V_p$ 上的线性映射；否则 $T_\gamma$ 无法定义。另注：原文在紧接本页之处写 “T_γ := Σ_{r∈I} T_{γ_k}”，其中下标 $k$ 应为 $r$——求和指标是 $r\in I$，且同段上文已写 $T_{\gamma_r}$；这是原文的笔误，正确形式为 $T_\gamma=\sum_{r\in I}T_{\gamma_r}$。]

于是我们得到一个映射 $\rho\colon\pi_1(S,p)\to\Aut(V_p)$，它是群同态。把 $\Aut(V_p)$ 与 $\GL_k(\C)$ 作一个等同之后，我们得到空间 $\mathcal{E}(S,\GL_k(\C))$ 中的一个元素。$\mathcal{E}^{\mathrm{tw}}(\Sigma,\C^\times)$ 的每个元素，连同覆叠 $\pi\colon\Sigma\to S$ 以及联系于 $\pi$ 的谱网络 $\mathcal{W}$，都给出空间 $\mathcal{E}(S,\GL_k(\C))$ 中的一个元素：
\[
\pi_1(S,p)\xrightarrow[\text{network}]{\text{spectral}}\mathbb{Z}^{\mathrm{tw}}_{B'}[\mathcal{P}(\Sigma)]\xrightarrow[\text{on }\Sigma]{\text{connection}}\GL_k(\C) .
\]

我们就这样构造出了一个映射 $\mathcal{E}^{\mathrm{tw}}(\Sigma,\C^\times)\to\mathcal{E}(S,\GL_k(\C))$。此外，所得到的 $\GL_k(\C)$-局部系统容许一个自然的标架（framing），由
\[
F(p):=\big(L_{p_1},\,L_{p_1}\oplus L_{p_2},\,\dots,\,L_{p_1}\oplus\cdots\oplus L_{p_k}\big)
\]
给出；于是我们得到映射 $\mathcal{E}^{\mathrm{tw}}(\Sigma,\C^\times)\to\mathcal{E}^{\mathrm{fr}}(S,\GL_k(\C))$。要说明这确实是一个非阿贝尔化映射，我们需要验证它覆盖 $\mathcal{E}^{\mathrm{fr}}(S,\GL_k(\C))$ 的一个开稠密子集。在下一节中，我们对于适应 $S$ 的理想三角剖分的谱网络证明这一点。

\section{非退化谱网络的阿贝尔化映射}

\subsection{秩二}

% [原文 p.171]
与以往一样，设 $S$ 是纤毛曲面。我们在 $S$ 上固定一个理想三角剖分 $\Delta$，并设 $\pi\colon\overline{\Sigma}\to\overline{S}$ 是由 $\Delta$ 连同谱网络 $\mathcal{W}$ 构造出的分歧 $2:1$ 覆叠。

我们取 $\mathcal{E}^{\mathrm{fr}}_\Delta(S,\GL_2(\C))$ 的一个元素，再次把它看作一个带标架的秩 $2$ 平坦向量丛 $(V,\nabla)$，即：对每个穿孔 $p$ 及其一个（可缩到 $p$ 的）小邻域 $U_p$，存在 $V$ 的一个 $\nabla$-不变线子丛 $\ell_p$。由 §10.2.2 中定义的 $\Delta$-横截性，对由 $\Delta$ 的一条边相连的两个穿孔，其 $\nabla$-不变线丛（沿该边搬运到某个公共点）处于直和之中。

\begin{remark}[15.1]
到 $\GL_2(\C)$ 的表示可能不容许任何 $\Delta$-横截的标架。例如，对于来自通过上三角矩阵群分解的表示的联络，这一性质不成立。不过，容许 $\Delta$-横截标架的联络是一般的（generic），即它们构成 $\mathcal{E}(S,\GL_2(\C))$ 的一个开稠密子集。
\end{remark}

现在取 $S\setminus\pi(\mathcal{W})$ 的一个连通分支 $C$，称为一个\textbf{胞腔}（cell，见图 15.1（右））。$C$ 是一个四边形，$C$ 的两个顶点是分支点，另外两个顶点是穿孔，设为 $p_1$ 与 $p_2$。由于 $C$ 单连通，我们可以取任意点 $z\in C$，并把 $\nabla$-不变子丛 $\ell_1:=\ell_{p_1}$ 与 $\ell_2:=\ell_{p_2}$ 从邻域 $U_{p_1}$ 与 $U_{p_2}$ 搬运到 $z$。由于 $C$ 单连通，这一平行移动不依赖于道路的选取。

\Fig[0.8\linewidth]{figures/fig_15_1.png}{15.1}{$\Sigma$ 中与 $S$ 中的胞腔。}

$C$ 关于 $\pi$ 的提升由 $C$ 的两个不交的拷贝 $C_1$ 与 $C_2$ 组成（见图 15.1（左））。我们假定 $C_1$ 含 $p_1$ 的提升，它是一个汇点（sink）；$C_2$ 含 $p_2$ 的提升，它是一个汇点。我们定义线丛 $L$：在 $C_1$ 上为 $\ell_1$，在 $C_2$ 上为 $\ell_2$。

为了整体地构造 $L$，我们需要沿谱网络的谱线把 $L$ 粘合起来，且粘合方式保持平行移动。为此，我们取两个共%
% [原文 p.172]
享 $\mathcal{W}$ 的一条边 $w$ 的胞腔 $C$ 与 $C'$。于是 $C$ 与 $C'$ 共享一个穿孔；不妨设它是 $p_2$。$C'$ 的第二个穿孔记作 $p'_2$。$C'$ 上相应的 $\nabla$-不变线丛记作 $\ell'_1$（对应于 $p_1$）与 $\ell'_2$（对应于 $p'_2$）。我们在 $C$ 上取 $\ell_1$ 与 $\ell_2$ 的截面 $s_1$ 与 $s_2$，在 $C'$ 上取 $\ell'_1$ 与 $\ell'_2$ 的截面 $s'_1$ 与 $s'_2$。我们可以假定 $s_1$ 在 $w$ 的一个邻域 $U_w$ 中与 $s'_1$ 一致，且 $s'_2(z)=s_2(z)+\alpha(z)s_1(z)$，其中 $\alpha\colon U_w\to\C$ 是某个函数。

现在设 $\nabla s_i=\de_is_i$、$\nabla s'_i=\de'_is_i$ 在 $U_w$ 中成立，其中 $\de_1,\de_2,\de'_1,\de'_2$ 是某些 $1$-形式。我们沿 $U_w$ 把 $L|_C$ 与 $L|_{C'}$ 粘合起来，把 $s_i$ 与 $s'_i$ 等同（$i\in\{1,2\}$）。这一粘合保持平行移动，当且仅当 $\de_1=\de'_1$ 且 $\de_2=\de'_2$。由于 $s_1=s'_1$ 在 $U_w$ 中成立，故 $\de_1=\de'_1$。进一步，一方面
\[
\nabla s'_2=\nabla(s_2+\alpha s_1)=\nabla s_2+\de\alpha\cdot s_1+\alpha\nabla s_1=\de_2s_2+\de\alpha\cdot s_1+\alpha \de_1s_1 .
\]
另一方面 $\nabla s'_2=\de'_2s'_2=\de'_2s_2+\de'_2\alpha s_1$。由于 $s_1$ 与 $s_2$ 线性无关，故 $\de_2=\de'_2$。因此，按这种方式我们可以构造 $\Sigma\setminus B'$ 上的一个平坦线丛，它在 $B'$ 的点周围单值为 $-1$，从而给出 $\mathcal{E}^{\mathrm{tw}}(\Sigma,\C^\times)$ 的一个元素。直接计算表明，这一过程是 §15.1 的非阿贝尔化过程的逆过程。
% [Tutor Note: 比较 $\nabla s'_2$ 的两种展开时，应分别比较 $s_2$ 与 $s_1$ 的系数：$s_2$ 的系数给出 $\de_2=\de'_2$（原文采用了这一步）；$s_1$ 的系数给出 $\de\alpha+\alpha \de_1=\de'_2\alpha$，即 $\de\alpha=\alpha(\de_2-\de_1)$，这正是联络 $1$-形式在标架变换 $s'_2=s_2+\alpha s_1$ 下的标准变换律（对 $\de_1=\de'_1$ 的情形）。核对方法：把 $\nabla$ 在标架 $(s_1,s_2)$ 下写成矩阵 $\de+\diag(\de_1,\de_2)$，在标架 $(s_1,s'_2)$ 下写成 $\de+\diag(\de_1,\de'_2)$，用 (3.4) 式比较即得同一结论。]

\subsection{高阶}

设 $S$ 是纤毛曲面，$\Delta$ 是 $S$ 的理想三角剖分。设 $k\ge2$。设 $\pi\colon\Sigma_k\to S$ 是 §13.5.2 中构造的分歧 $k$ 重覆叠。设 $(V,\nabla,F)\in\mathcal{E}^{\mathrm{fr}}_\Delta(S,\GL_k(\C))$。我们再次有一个 $(V,\nabla)$ 的标架 $F$，即在每个穿孔 $p$ 的小邻域中，$V$ 的平坦子丛的一个嵌套序列
\[
F^p_0\subset F^p_1\subset\cdots\subset F^p_n ,
\]
其中对一切 $1\le i\le n-1$ 有 $\dim F^p_i=i$。

设 $t$ 是 $\Delta$ 的一个三角形，$(p,q,r)$ 是 $t$ 的各顶点处的穿孔。由于 $t$ 单连通，我们可以用 $\nabla$ 把三个旗子丛 $F^p$、$F^q$ 与 $F^r$ 搬运到 $t$ 中任一公共点。回忆一下：由于 $\nabla$ 是 $\Delta$-横截的，对任意三个满足 $a+b+c=k$ 的非负整数 $a,b,c$，都有
\[
F^p_a\oplus F^q_b\oplus F^r_c=V .
\]
特别地，这一假设保证：对任意三个满足 $a+b+c=2k+1$ 的非负整数 $a,b,c$，从 $F^p_a\cap F^q_b\cap F^r_c$ 到 $F^p_a/F^p_{a-1}$（相应地 $F^q_b/F^q_{b-1}$ 与 $F^r_c/F^r_{c-1}$）的典范投影是同构。

\begin{remark}[15.2]
容许 $\Delta$-横截标架的联络所成的子集在 $\mathcal{E}(S,\GL_k(\C))$ 中是开稠密的。
\end{remark}

我们如下在 $\Sigma_k$ 上构造线丛 $L$。设 $p$ 是 $S$ 的一个穿孔，$1\le a\le n$。设 $S_{\gamma_p^{(a)}}$ 是引理 13.17 中定义的穿孔圆盘。在这个穿孔圆盘上，线丛 $L$ 是平凡丛 $F^p_a/F^p_{a-1}$。设 $e$ 是 $\Gamma_k$ 的一条边，$t=(p,q,r)$ 是含 $e$ 的三角形。沿 $e$ 恰好粘合着 $\Sigma_k$ 中的两个穿孔圆盘。设 $S_{\gamma_p^{(a)}}$ 与 $S_{\gamma_q^{(b)}}$ 是这两个圆盘。在这些圆盘内部，线丛 $L$ 分别是 $F^p_a/F^p_{a-1}$ 与 $F^q_b/F^q_{b-1}$。它们沿 $e$ 通过同构
% [原文 p.173]
\[
F^p_a/F^p_{a-1}\cong F^p_a\cap F^q_b\cap F^r_{2k+1-a-b}\cong F^q_b/F^q_{b-1}
\]
粘合起来。$L$ 的这一构造具有局部常值的转移函数，故它定义了 $\Sigma_k$ 去掉 $\Gamma_k$ 的顶点之后的一个平坦线丛 $L$。一个简单的线性代数计算表明：$L$ 绕 $\Gamma_k$ 的白顶点的单值是平凡的，而绕内部黑顶点（即 $\pi\colon\Sigma_k\to S$ 的分歧轨迹）的单值是 $-1$。我们就这样构造出了 $\mathcal{E}^{\mathrm{tw}}(\Sigma,\C^\times)$ 的一个元素。再次，直接计算表明 $L$ 的非阿贝尔化是 $(V,\nabla)$。

\section{谱网络与 Fock–Goncharov 坐标}

设 $S$ 是纤毛曲面，$\Delta$ 是 $S$ 的理想三角剖分。非阿贝尔化映射与阿贝尔化映射定义了 $S$ 上秩 $k$ 的 $\Delta$-横截带标架平坦向量丛的模空间 $\mathcal{E}^{\mathrm{fr}}_\Delta(S,\GL_k(\C))$ 与 $\Sigma_k$ 上在分歧点周围单值为 $-1$ 的平坦线丛的模空间 $\mathcal{E}^{\mathrm{tw}}(\Sigma_k,\C^\times)$ 之间的一个同胚。后一个模空间要容易研究得多：它没有标架（因而没有横截性条件），并且线丛的空间容许一个非常方便的参量化，即作为第一同调群 $H_1(\Sigma_k)$ 的特征所成的集合（见注 3.31 与例 13.1）。特别地，空间 $\mathcal{E}(\Sigma_k,\C^\times)$ 同胚于 $(\C^\times)^n$，其中 $n$ 是 $H_1(\Sigma_k)$ 的维数。这一同胚由 $H_1(\Sigma_k)$ 的一组基的选取给出，而我们可以把这组基取为 $\Sigma_k$ 上 $n$ 条曲线的一个集合。这 $n$ 条曲线的和乐的选取随即唯一确定 $\Sigma_k$ 上的一个平坦线丛，从而通过非阿贝尔化映射确定 $S$ 上一个秩 $k$ 的 $\Delta$-横截带标架平坦向量丛。这并不是我们第一次参量化 $\mathcal{E}^{\mathrm{fr}}_\Delta(S,\GL_k(\C))$；在 §10.2.3 中我们已经用 Fock–Goncharov 坐标描述过这个空间的一种参量化。本节的目标是要说明：对谱曲面 $\Sigma_k$ 与 $H_1(\Sigma_k)$ 的基作适当选取之后，这两种参量化是相同的：每个 Fock–Goncharov 坐标对应于 $\Sigma_k$ 上某条特定曲线的和乐，而所有这些曲线的集合构成 $H_1(\Sigma_k)$ 的一组基。

谱网络与 Fock–Goncharov 坐标之间的关系由 Gaiotto、Moore 与 Neitzke 在 [156] 中作了研究。在该文中，他们利用底曲面 $S$ 上一个经过恰当选取的谱网络之补集的连通分支，与联系于带标架表示 $(V,\nabla,F)$ 的 snakes（见 §10.2.4 的讨论）之间的一个对应，来说明：阿贝尔化局部系统绕谱覆盖上特定曲线的单值正是 $V$ 的 $\mathcal{X}$-坐标。我们这里采取的方法略有不同，这源于我们构造阿贝尔化映射的方式。

我们首先关注三个横截旗的情形，即 $S$ 是带有穿孔 $p,q,r$ 的理想三角形。设 $n\ge2$，$(A,B,C)$ 是 $\C^n$ 中三个横截旗。如 §10.2.3 中所讨论，这三个旗由 $\frac{(n-1)(n-2)}{2}$ 个三重比的数据所决定，这些三重比称为该三元组的 $\mathcal{X}$-坐标。在 §13.5.2 中，我们构造了 $S$ 的 $k$ 重谱覆盖 $\Sigma_k$ 连同
% [原文 p.174]
$\Sigma_k$ 上的谱网络 $\mathcal{W}$。以三个旗 $(A,B,C)$ 为标架的 $S$ 上平凡局部系统于是可以被阿贝尔化为 $\Sigma_k$ 上的平坦线丛 $L$。当 $n$ 为偶数时曲面 $\Sigma_k$ 有 $\frac{n}{2}$ 个边界分支，当 $n$ 为奇数时有 $\frac{n+1}{2}$ 个边界分支；而 $\Sigma_k$ 上有 $\frac{n(n-1)}{2}$ 个简单分歧点。由 Riemann–Hurwitz 公式，当 $n$ 为偶数时 $\Sigma_k$ 的亏格为 $\frac{(n-2)^2}{4}$，当 $n$ 为奇数时为 $\frac{(n-1)(n-3)}{4}$。

在两种情形下，第一同调群 $H_1(\Sigma_k)$ 的维数（它同时也是 $\Sigma_k$ 上平坦线丛的模空间的维数）都是 $\frac{(n-1)(n-2)}{2}$。沿用 §13.5.2 的记号，设 $a,b,c\ge0$ 满足 $a+b+c=2k+3$，并设 $v_{a-1,b,c}$（相应地 $v_{a,b-1,c}$、$v_{a,b,c-1}$）是 $\Gamma_k$ 中位于交集 $\gamma_p^{(a-1)}\cap\gamma_q^{(b)}\cap\gamma_r^{(c)}$（相应地 $\gamma_p^{(a)}\cap\gamma_q^{(b-1)}\cap\gamma_r^{(c)}$、$\gamma_p^{(a)}\cap\gamma_q^{(b)}\cap\gamma_r^{(c-1)}$）上的黑顶点。设 $\tilde{x}_{a,b,c}$ 是 $S$ 上绕这三个分歧点的回路（见图 15.2）。
% [Tutor Note: 这里的维数 $\frac{(n-1)(n-2)}{2}$ 与前一段“三个横截旗由 $\frac{(n-1)(n-2)}{2}$ 个三重比决定”中的数目完全一致，这正是本节论证的关键一致性：$\mathcal{X}$-坐标的个数必须等于 $H_1(\Sigma_k)$ 的秩，否则“每个 $\mathcal{X}$-坐标对应一条基曲线的和乐”这一陈述无法成立。核对方法：取 $n=2$ 得 $0$，与 $\Sigma_2$ 上只有平凡同调一致（此时三个旗的模空间是零维的）；取 $n=3$ 得 $1$，此时恰有一个三重比。]

\Fig[0.5\linewidth]{figures/fig_15_2.png}{15.2}{联系于 $\Gamma_4$ 的道路 $\tilde{x}_{3,4,4}$。分歧点（即内部黑顶点）用蓝色叉号表示。}

把道路 $\tilde{x}_{a,b,c}$ 提升到 $\Sigma_k$（用标准的路径提升法则，而非谱网络法则），我们得到 $k$ 个不同的提升，其中 $k-2$ 个是 $\Sigma_k$ 中的闭回路，另外两个不是回路，但衔接起来之后构成一个回路。我们用 $x_{a,b,c}$ 记衔接这两个提升所得的回路。下述命题可由阿贝尔化局部系统 $L$ 的定义与三重比的定义推出：

\begin{proposition}[15.3]
$L$ 沿 $x_{a,b,c}$ 的单值等于三重比 $r_{a,b,c}(F(p),F(q),F(r))$。
\end{proposition}

现在考虑共享一条以 $p$ 与 $q$ 为端点的边的两个三角形 $t_1$ 与 $t_2$。设 $r$ 与 $s$ 是其余两个穿孔，使得 $t_1=(p,q,r)$ 且 $t_2=(p,s,q)$，并设 $(F(p),F(s),F(q),F(r))$ 是 $S=t_1\cup t_2$ 上平凡局部系统的一个横截标架。设 $\Gamma_k$ 是 §13.5.2 中在 $S$ 上构造的图。对满足 $a+b=k+2$ 的 $a,b\ge0$，设 $\tilde{x}_{a,b}$ 是绕两个内部黑顶点 $v^{(t_1)}_{a,b,n}$ 与 $v^{(t_2)}_{a,n,b}$ 的回路，如图 15.3 所示。

再次，$\tilde{x}_{a,b}$ 到 $\Sigma_k$ 有 $k$ 个标准提升，其中恰好两个不是可缩回路。我们在 $\Sigma_k$ 上选取其中一个提升 $x_{a,b}$。由我们对 $\Sigma_k$ 上阿贝尔化局部系统 $L$ 的构造，我们有下述命题：

\Fig[0.65\linewidth]{figures/fig_15_3.png}{15.3}{联系于 $\Gamma_4$ 的道路 $\tilde{x}_{2,4}$。分歧点（即内部黑顶点）用蓝色叉号表示。}

% [原文 p.175]
\begin{proposition}[15.4]
$L$ 沿 $x_{a,b}$ 的单值等于交比 $r_{a,b}(F(p),F(s),F(q),F(r))$。
\end{proposition}

通过对谱覆盖 $\Sigma_k$ 的拓扑作仔细研究，可以证明：取遍所有可能的回路 $x_{a,b,c}$ 与 $x_{a,b}$ 所得的族构成 $H_1(\Sigma_k)$ 的一组基。因此，Fock–Goncharov 坐标恰好就是绕谱曲面 $\Sigma_k$ 的第一同调群的一组基的单值。

\begin{remark}[15.5]
这一对应是就谱曲面 $\Sigma_k$ 的这一特定选取而言的。一般说来，把 Fock–Goncharov 坐标与在另一个谱曲面上的阿贝尔化的和乐联系起来并不容易。其他簇变量是否容许类似的几何描述，目前也尚不清楚。
\end{remark}

\begin{remark}[15.6]
在 [156] 中，Gaiotto、Moore 与 Neitzke 证明了：当 $k\le5$ 时，谱曲面 $\Sigma_k$ 及其谱网络按 §13.6 的意义是 WKB 谱网络。这一结论猜想对 $k>5$ 仍然成立。他们还证明，这些谱网络可以通过对 Argyres–Douglas 理论 $AD_1$ 的“第 $k$ 层提升”（level $k$ lift）作扰动而得到。
\end{remark}

\section{Fenchel–Nielsen 谱网络的阿贝尔化映射}

我们将构造联系于闭曲面 $S$ 上一个 Fenchel–Nielsen 谱网络 $\mathcal{W}$（带选定的美式或英式消解）的阿贝尔化映射。我们采用与 §15.2.1 中秩二非退化网络所给出的方法非常类似的做法。与非退化情形一样，除了 $S$ 上平坦 $\GL_2$-丛的数据之外，我们还需要一些额外数据，我们像对非闭曲面那样称之为标架。为此我们注意：$\mathcal{W}$ 在 $S$ 中的补集是有限多个连通分支的不交并，每个分支都同胚于圆环；而其中任一分支 $C$ 到 $\Sigma$ 的提升，是两个圆环 $C_1,C_2\subset\Sigma$ 的不交并。
% [Tutor Note: 这里“每个分支同胚于圆环”是 Fenchel–Nielsen 网络的标志性特征，可与 §14.2 对照核对：非退化网络的补集分支是胞腔（四边形），因为谱线终止于穿孔；而 FN 网络的谱线全是双墙、没有穿孔端点，故补集分支沿对偶方向是无孔的，只能是圆环。核对方法：由 §14.2，$S\setminus\mathcal{W}(\phi)$ 的每个分支是被闭 $\vartheta$-轨道扫出的穿孔圆盘；而在闭曲面情形（$P=\varnothing$ 的退化极限）这些圆盘两两粘合成圆环。]

% [原文 p.176]
我们利用谱网络的消解来对 $C$ 的定向作一个选取：$C_1$ 或 $C_2$ 中之一的边界由谱线的前半段构成，即某些谱线 $p_h\in\mathcal{W}$ 的像 $p_h\big((-1,0]\big)$。至多重新命名之后，我们把该提升称为 $C_1$，并按谱线在 $C_1$ 边界上诱导的定向来定向 $C$。我们称 $C$ 的这一方向为\textbf{正}（或 $+$）向，而相反方向为\textbf{负}（或 $-$）向。

\begin{definition}[15.7]
设 $V\in\mathcal{E}(S,\GL_2(\C))$ 是 $S$ 上的平坦 $\GL_2(\C)$-丛。$V$ 的一个\textbf{标架}是对每个这样的定向圆环 $C$，选取 $V$ 绕 $C$ 沿 $+$ 方向的单值 $M_+$ 的一条特征线 $\ell_+(C)$，以及沿 $-$ 方向的单值 $M_-=M_+^{-1}$ 的一条特征线 $\ell_-(C)$。我们要求 $\ell_+\neq\ell_-$，并且联系于一条双墙两侧的标架是横截的。
\end{definition}

\begin{remark}[15.8]
为使平坦 $\GL_2(\C)$-丛容许一个标架，绕每个圆环的单值必须是可对角化的。当相应的表示是 Fuchs 型时尤其如此，因为此时绕每条闭曲线的单值都是 $\SL_2(\C)$ 中的双曲元素。
\end{remark}
% [Tutor Note: $\ell_-(C)$ 取作 $M_-=M_+^{-1}$ 的特征线而非 $M_+$ 的另一条特征线，是必须的：$C$ 的 $-$ 向与 $+$ 向是同一条圆环上方向相反的两条曲线，故其单值互为逆（同伦类取逆）。核对方法：若 $M_+$ 有特征值 $\lambda,\lambda^{-1}$（$\GL_2$ 情形下 $\det M_+$ 未必为 $1$，则记特征值为 $\lambda,\mu$），则 $M_-=M_+^{-1}$ 的特征值为 $\lambda^{-1},\mu^{-1}$，特征线相同但标注互换；因此“$\ell_+\neq\ell_-$ 且二者横截”等价于“$M_+$ 可对角化且不是标量矩阵”。当 $M_+$ 为标量时两条特征线重合，标架不存在——这与注 15.8 的“必须可对角化”相衔接。]

设 $V\in\mathcal{E}(S,\GL_2(\C))$，$\ell$ 是 $V$ 的一个标架。我们想在 $\Sigma$ 上构造一个平坦线丛 $L$。我们先在 $\mathcal{W}$ 的补集上构造这个丛。$S\setminus\mathcal{W}$ 的每个连通分支 $C$ 有两个提升 $C_1$ 与 $C_2$，按上述方式标注。我们定义 $L$ 为如下的（平坦）平凡丛：在 $C_1$ 上等于 $\ell_+(C)$，在 $C_2$ 上等于 $\ell_-(C)$。现在我们要沿谱网络的谱线把刚构造出的平坦线丛粘合起来，这些谱线全是双墙。设 $C$、$C'$ 是被一条双墙隔开的 $S\setminus\mathcal{W}$ 的两个连通分支。双墙在 $\Sigma$ 上的两个提升把 $C_1$ 与 $C'_2$（相应地 $C_2$ 与 $C'_1$）分开，我们用如下同构把 $\ell_+(C)$ 粘到 $\ell_-(C')$：
\[
\begin{array}{ccccc}
\ell_+(C) & \to & \ell_+(C')\oplus\ell_-(C') & \to & \ell_-(C')\\[2pt]
x & \mapsto & x=x'_+ + x'_- & \mapsto & x'_- .
\end{array}
\]
我们类似地把 $\ell_-(C)$ 粘到 $\ell_+(C')$。这些粘合映射是局部常值的，故所得的线丛 $L\to\Sigma$ 带有平坦联络，它在各圆环的提升上延拓了平凡联络。绕分歧点的单值为 $-1$，于是我们得到 $\mathcal{E}^{\mathrm{tw}}(\Sigma,\C^\times)$ 的一个元素，我们称它为 $V$ 的\textbf{阿贝尔化}。

\section{Fenchel–Nielsen 谱网络的非阿贝尔化映射}

我们以与用于非退化网络时略有不同的方式定义 Fenchel–Nielsen 谱网络的非阿贝尔化映射。这是因为适应于 Fenchel–Nielsen 谱网络的路径提升法则一般会产生无穷多个提升，从而要求我们赋予谱覆盖上道路的无穷形式和以意义。愿意了解这一路径提升途径的读者%
% [原文 p.177]
可以在 [210] 中读到细节。我们这里选择一种不同的途径，不仅是为了简单，也是为了展示完成这一构造的另一种方式：解孤子内容（solving for soliton content）。

设 $S$ 是闭曲面，$\mathcal{W}$ 是 $S$ 上的 Fenchel–Nielsen 谱网络（连同它的一个消解），$\pi\colon\Sigma\to S$ 是谱覆盖。设 $L\in\mathcal{E}^{\mathrm{tw}}(\Sigma,\C^\times)$ 是 $\Sigma$ 上的平坦线丛。我们选取一条分支割线并对 $\Sigma$ 的两个叶编号，于是谱网络的每条谱线都由一个有序对 $(1,2)$ 或 $(2,1)$ 标注，它由谱线起始的叶以及终止的叶给出；我们把 $S$ 上的谱线定向，使它们背离所经过的分支点。我们首先把线丛 $L$ 前推（push forward）到 $S\setminus\mathcal{W}$，得到 $S\setminus\mathcal{W}$ 上一个秩 $2$ 的平坦向量丛 $\pi_\ast L$。现在我们要对 $\mathcal{W}$ 中每条标号为 $(i,j)$ 的谱线 $p_h$ 联系一个映射 $e_h\colon L(z_i)\to L(z'_j)$，其中 $z$（相应地 $z'$）是 $p_h$ 右侧（相应地左侧）一个靠近 $p_h$ 的点。穿过 $S$ 上一条双墙便将诱导出如下形式的映射：
\[
\begin{pmatrix}1 & e_h\\ 0 & 1\end{pmatrix}
\begin{pmatrix}1 & 0\\ -e_{h'} & 1\end{pmatrix} .
\]

现在再加上一个约束：秩 $2$ 丛在 $S$ 上绕任一分支点的单值必须是平凡的。对一个 Fenchel–Nielsen 谱网络的每个基本片段（分子 I 与分子 II，见图 14.3），这给出两个 $2\times2$ 矩阵值方程，其中含 $6$ 个待定变量，即联系于一个基本片段中 $6$ 条谱线的 $6$ 个映射 $e_h$。这些方程容许唯一解。读者可在 [205] 中找到这些解的详细计算。沿每条双墙所联系的映射粘合丛 $\pi_\ast L$，便得到 $S$ 上一个平坦秩 $2$ 丛 $V$。由于这个丛限制在 $S\setminus\mathcal{W}$ 的连通分支 $C$ 上是平凡丛 $\pi_\ast L|_{C_1\cup C_2}$，它自然地劈分为两个平坦线子丛，且二者都被绕 $C$ 的单值所保持。这说明所得的丛 $V$ 自然地带有标架。我们就这样构造出了 $\mathcal{E}(S,\GL_2(\C))$ 的一个元素，我们称它为 $L$ 的\textbf{非阿贝尔化}。

\section{谱网络与 Fenchel–Nielsen 坐标}

对一个 Fenchel–Nielsen 谱网络，我们下面将看到如何在平坦 $\SL_2(\C)$-联络的模空间上得到复化的 Fenchel–Nielsen 坐标，这里遵循 [205, Section 8.4]。

回忆一下，在 §15.3 中我们用了纤毛曲面 $S$（带理想三角剖分 $\Delta$）的 $k$ 重谱覆盖以及 $\Sigma_k$ 上的谱网络，以绕 $H_1(\Sigma_k)$ 一组基的单值的形式，得到空间 $\mathcal{E}^{\mathrm{fr}}_\Delta(S,\GL_k(\C))$ 上的 Fock–Goncharov 坐标。

这里我们取闭曲面 $S$ 与 $S$ 上一个 Fenchel–Nielsen 谱网络 $\mathcal{W}$（带选定的美式或英式消解）。考虑带标架的一点 $\nabla$，它位于平坦 $\GL_2(\C)$-丛%
% [原文 p.178]
的空间中。§15.4 的阿贝尔化过程给出 $\mathcal{E}^{\mathrm{tw}}(\Sigma,\C^\times)$ 的一个元素 $\nabla^{\mathrm{ab}}$。沿一条绕圆环正向行走的道路的提升，$\nabla^{\mathrm{ab}}$ 的和乐等于 $M_+$ 的单值特征值。因此，这个单值特征值是 $\nabla$ 的指数化长度坐标的平方根的一个复化版本。注意，由阿贝尔化映射可见，给定一个 Fenchel–Nielsen 谱网络，这样的 Fenchel–Nielsen 长度坐标是典范地得到的。然而，它们只给出 $H_1(\Sigma_2)$ 的一组基的一半（这与 Fock–Goncharov 坐标的情形不同：那里所有相关曲线的集合给出了 $H_1(\Sigma_k)$ 的一组基）。换言之，仅凭一个 Fenchel–Nielsen 谱网络的选取以及一个裤子分解，并不足以给出平坦 $\PSL_2(\C)$-联络的模空间上的一个完整坐标系。

另一半基元素不能典范地得到，它们来自在联络 $\nabla$ 的某种修改之下考虑 $\nabla^{\mathrm{ab}}$ 沿一个圈 $\gamma$ 的和乐，这种修改称为\textbf{扭转流}（twist flow）。这一修改出现在下述操作中：沿一个圆环切开曲面 $S$，然后用在保持单值特征空间的自同构，把边界分支上的联络重新粘回去。注意，在我们于 §15.4 中得到阿贝尔化的方式中，选取的是平凡的这种自同构。

因此，在这种切开与粘合的方式下，原来的联络 $\nabla$ 总能被阿贝尔化为一个联络 $\nabla^{\mathrm{ab}}$，但平行移动会受到上述自同构的不确定性的影响。事实上，人们得到一个单参数族的修改后联络 $\nabla(\lambda)$，而用丛在圆环上的截面 $s_1,s_2$ 来表示，这样一个自同构的作用可以写成
\[
s_1\to\lambda s_1,\qquad s_2\to\lambda^{-1}s_2 .
\]
对如图 15.4 中的圈 $\gamma\in H_1(\Sigma_2)$，可以直接验证：扭转流对 $\nabla^{\mathrm{ab}}$ 的和乐 $\mathcal{X}_\gamma$ 的作用为
\[
\mathcal{X}_\gamma\mapsto\lambda^2\mathcal{X}_\gamma .
\]
扭转流下的这一变换律给出一个复化的指数化 Fenchel–Nielsen 扭转坐标。最后，在保持网络 $\mathcal{W}$ 不变的情况下改变圈 $\gamma$ 的选取，会把扭转坐标乘以指数化长度坐标的某个幂。导致这类长度–扭转坐标的 Fenchel–Nielsen 谱网络的详细例子可在 [205, Sect. 9.3] 中找到。
% [Tutor Note: $\lambda^2$ 中指数 $2$ 的来源可这样核对：$\gamma$ 在 $S$ 中穿过圆环一次，而它在 $\Sigma_2$ 上的提升从一个叶出发、在另一个叶回到原叶，故在提升的道路上要跨越双墙两次，每次跨越处的粘合按扭转变换贡献一个因子 $\lambda$（或 $\lambda^{-1}$，取决于从哪个叶到哪个叶），两次相乘得 $\lambda^{\pm2}$。核对方法：把 $\nabla(\lambda)$ 限制在圆环上，其单值矩阵为 $\diag(\lambda,\lambda^{-1})M$；由于 $\mathcal{X}_\gamma$ 是 $L$ 的纤维之间的标量，在标架 $s_1\mapsto\lambda s_1$、$s_2\mapsto\lambda^{-1}s_2$ 下它按 $\lambda^2$ 变换，与文中公式一致。]

\Fig[0.5\linewidth]{figures/fig_15_4.png}{15.4}{圈 $\gamma$，其提升的单值代表扭转参数。}

% [原文 p.179]
\begin{remark}[15.9]
事实上，上述阿贝尔化 $\nabla^{\mathrm{ab}}$ 的和乐所含的信息比 §7.1 中的 Fenchel–Nielsen 坐标略多：因为此处通过阿贝尔化过程描述的长度–扭转坐标给出相差等价意义下的一个完整 $\SL_2(\C)$-联络，而复化的 Fenchel–Nielsen 坐标只确定它在 $\PSL_2(\C)$ 中的投影。
\end{remark}
% [Tutor Note: 这两种坐标的差别正是中心 $\Z/2=\{\pm I\}$ 的提升：$\SL_2(\C)\to\PSL_2(\C)$ 的核为 $\{\pm I\}$，故一个 $\PSL_2(\C)$-联络在 $\SL_2(\C)$ 中有两个提升，二者相差对各和乐元素同时变号。这也解释了正文为何说单值特征值是“指数化长度坐标的平方根”的复化版本：长度坐标取的是 $\lambda^2$（平方），而在 $\SL_2$ 的提升中记录的是 $\lambda$（即平方根），符号的选取对应于分歧点周围单值 $-1$ 所对应的自旋结构。核对方法：取 $M_+$ 的特征值 $\lambda$，则指数化长度坐标 $=\lambda^2$ 在 $\lambda\mapsto-\lambda$ 下不变，而 $\nabla^{\mathrm{ab}}$ 的和乐随 $\lambda\mapsto-\lambda$ 变号，二者恰好差一个 $\Z/2$。]

\begin{remark}[15.10]
研究（对上述 Strebel 微分 $\phi$ 的）Fenchel–Nielsen 谱网络 $\mathcal{W}(\phi)$ 的一个有力物理动机在于：相应的复化 Fenchel–Nielsen 坐标与相关 $\mathcal{N}=2$ 超对称场论的库仑分支内部一个重要的“实”子空间相联系。这为从颇为几何的视角理解 AGT 对应提供了新的途径；关于这一视角的更多信息，我们参见 [205, Sect. 1.2] 及其中的文献。
\end{remark}

% [原文 p.180]
% 本页为原书空白页（第 15 章章末的空页），无正文内容。
